\chapter{Does Language Shape Our Worldview?}
\section{A Box of Coloring Pens}
First pick up a box of coloring pens.

Not an art student's expensive set, just the familiar stationery-shop kind: twelve pens in a plastic case, caps labeled red, orange, yellow, green, blue, purple. You have used them for classroom posters and cell diagrams in science, so often that you would never question them.

Try a little experiment. Put this twelve-color box beside the forty-eight-color set on the next shelf, then look out at the sky for a minute. Its color changes continuously from overhead to the horizon: deeper above, paler near the buildings, different again at the clouds' edges. Here is the question: does the child with twelve colors see the same sky as the child with forty-eight?

Intuition says: of course. Eyes are eyes, light is light; the sky does not ration its colors because your pens were cheaper. Colors exist, and words merely label them. More or fewer labels, the goods remain the same.

But do not close the box yet. Ask further: suppose your language has no word for ``blue.'' Not that you forgot to learn it, but that the language never gave blue a separate name, just as ours gives no separate basic names to light gray-blue and dark gray-blue. Is the sky still blue? When you and a friend say ``What a lovely day,'' is the color in your minds really identical?

It sounds like a word game. Yet it has troubled philosophers, linguists, and psychologists for over a century, producing hypotheses, myths, experiments, and fierce disputes that continue today. Its formal question is: \textbf{does language shape our worldview?}

Starting with this box, we will take the question apart layer by layer. You will see that the answer is neither ``yes'' nor ``no,'' but something more interesting than either.

\section{Cook's Kangaroo and a People Without ``Left'' and ``Right''}
Come with me to a scene.

Time: afternoon. Place: near Cooktown, Queensland, in northeastern Australia. The tropical sun is white-bright; eucalyptus scents the air. A red dirt road crosses savanna. Termite mounds stand in the distance like red chimneys.

In 1770, British explorer James Cook's Endeavour struck the Great Barrier Reef and was forced into this river mouth for seven weeks of repairs. A naturalist aboard recorded a local Indigenous animal name in the voyage journal: a large hopping creature with short forelegs and long hind legs. The word traveled to Britain, then around the world. Today everyone knows it: kangaroo. Its home is here, in the local language Guugu Yimithirr. Incidentally, linguists later returned to check and found that the word originally named only a particular kind of large kangaroo. Outsiders expanded it to all kangaroos. Words change meaning when they travel; remember this clue.

Two centuries later, linguists conducting long-term research on the same red soil found something even more interesting: this language has no ``left'' and ``right.''

This does not mean speakers cannot distinguish their two differently functioning hands. Rather, they do not use expressions like ``the person on your left'' or ``move the cup right.'' They give directions through north, south, east, and west. A speaker might warn, ``There's an ant on your northern knee,'' or offer water by saying ``The cup is a little southeast of you.'' Long-term researchers recorded almost legendary details: someone casually retelling a film says ``The policeman stood west of the tree''; someone recalling an event years later reports which way each person faced. Without directions, their language cannot make the account complete.

Pause and feel what that means. In our language, you can say ``I met Xiaoming at the classroom door yesterday'' without considering which way he faced. In that language, before finishing, you must establish everyone's position. This is not fussiness but a grammatical requirement. Direction is to them what tense is to us: can you ignore time? Your verbs will not allow it. Saying ``have eaten'' and ``am eating,'' we are forced to mark each action's time; they must mark every thing's direction in every sentence.

Now the real question can appear.

\section{Do Eyes See, or Does Language See for Them?}
Lay out the material, and the conflict emerges.

On one side is solid common sense: the world exists, light enters the eye, electrical signals enter the brain. This works the same way everywhere on Earth. A Guugu Yimithirr-speaking child and a Chinese-speaking child have no difference in retinal structure. Direction and color perception are human hardware, not software installed by language.

On the other side are increasingly hard-to-ignore facts: the world's six thousand-plus languages require speakers to notice remarkably different things. Some, like ours, require a time tag for each action; others, a direction tag for each object. Still others, such as Turkish and Tuyuca in South America's Amazon basin, require you to state how you know each thing. Before saying ``Xiaoming broke the cup,'' you must choose a grammatical marker announcing whether you saw it, heard it reported, or inferred it. We can utter the claim without marking a source; they cannot. Grammar pursues them for an answer.

A genuine question therefore appears:

\textbf{Can your native language decide what you notice, what you remember, even what you are able to think?}

There are two versions. Keep them separate: this chapter's battles are fought between them.

\textbf{Strong version: language determines thought.} Your language determines the thought-world you inhabit; its limits are your world's limits. Some ideas are unavailable to speakers of another language, not because they will not think them, but because they cannot. This version usually appears under two names, American linguist Edward Sapir and his student Benjamin Lee Whorf, hence the ``Sapir--Whorf hypothesis,'' also called linguistic relativity.

\textbf{Weak version: language influences thought.} Language builds roads, not walls. It makes some things easier to notice, remember, and say, while others take more effort. Language is habit, not a prison.

The versions sound close but describe entirely different worlds. If the strong version is right, translation is impossible and learning another language cannot save you: you merely carry a Chinese prison into an English prison. If the weak version is right, differences between languages are scenery you can visit, not an uncrossable wall.

Do not choose sides yet. Hear both out first.

\section{Two Voices Worth Hearing in Full}
\textbf{Voice one: language is thought's mold.}

First give the strong version its full case. Today it is often ridiculed as ``a foolish idea disproved long ago.'' That is unfair and teaches you little.

Whorf himself is worth knowing. His regular occupation was not university professor but fire-insurance investigator. Linguistics consumed his spare time. His day job gave him a distinctive perspective. Investigating factory fires, he noticed a recurring pattern: workers carefully handled ``full gasoline drums'' but casually tossed cigarette ends near ``empty gasoline drums.'' Empty suggested ``nothing inside, no danger.'' In reality, the empty drums held gasoline vapor and could explode more readily than full ones. Whorf believed what killed these people was vocabulary, not chemistry: they acted according to the word's meaning rather than the thing's reality.

Following this thought, he studied Hopi, an Indigenous North American language, and reached a startling conclusion: it lacked grammatical machinery corresponding to our ``time.'' In Hopi people's world, time was not a river flowing from past to future. The strong version now had its full three-layer argument. First, grammar operates automatically. You cannot switch off tense or direction markers while speaking; they act faster than awareness. Second, grammar is installed before thought: children learn to speak, then use speech to think. Language is thought's operating system, not an application. Third, human grammars differ enormously. If none of these differences affects thought, are thousands of grammars merely thousands of meaningless coincidences? Even God's dice should not be so dull.

This is no weak case. It explains real ``empty-drum'' accidents and why our own categories seem natural. It also carries a moving humility: perhaps language has always been steering you while you thought you were steering it.

\textbf{Voice two: words are labels; humanity shares the foundation beneath the eyes.}

The other camp has solid material too. First, translation is difficult but keeps happening. If thought were truly locked by language, the Analects could have no English version and the United Nations could not function for a day. Translation loses things, but flavor, not an entire universe. Second comes experimental evidence. In 1969, anthropologist Brent Berlin and linguist Paul Kay surveyed color terms in twenty languages and checked published material on dozens more. Color terms did not grow at random. A language with only two basic terms distinguished ``dark'' and ``light''; with three, the third was always ``red.'' Yellow, green, and blue came later. Unrelated languages around the world developed color vocabulary in a strikingly similar order. This looks less like separate worlds than different houses on a shared visual foundation. Third is the simplest reason: no word does not mean no perception. Chinese has no single term for ``the color in the corner of your left eye,'' yet you see it. English has no separate basic words for one's mother's brother and father's brother, yet English speakers know which relative is which.

While listening to the dispute, expose a false witness repeatedly called for the strong version. Many people have heard that ``Eskimo languages have fifty or even hundreds of words for snow, so their snow-world differs completely from ours.'' This example appears in countless books, talks, and lessons, its number constantly rising. Linguist Geoffrey Pullum's famous essay The Great Eskimo Vocabulary Hoax dissected the snowball. Its origin was a few distinct roots mentioned by anthropologist Franz Boas early in the twentieth century, inflated through successive popular accounts. Besides, count roots and compounds generously and English snow, sleet, slush, blizzard, flurry, and snowfall form a substantial list too. This is a warning: a good part of the strong version's popular prestige rests on repeated misinformation. Hear ``this language has so many words for X, therefore these people are Y,'' and first check the number, not marvel at the people.

\section{Thinking Bridge: Attention Is Not a Cage}
Here is a tool to take away. Next time you hear ``speakers of this language are naturally such-and-such,'' whether praise or insult, dismantle it with three questions.

\textbf{First: what does this grammar force speakers to notice?}

Each language has a ``compulsory attention checklist.'' English speakers must state time, past or present, and number, one or many. Guugu Yimithirr speakers must state direction; Turkish speakers must state information source. These are not matters of refinement but taxes grammar collects whenever you open your mouth.

\textbf{Second: can outsiders learn this habit, and can speakers set it aside?}

This separates a real cage from an apparent one. A Turkish speaker learning another language can learn to speak without source marking. A Chinese speaker can train the habit of continually tracking direction; scouts and wilderness guides do. Learnable and dispensable means habit. Only what cannot be learned or set aside deserves to be called a cage.

\textbf{Third: at what level does the difference occur?}

Separate at least three layers: what the eyes see, perception; what the brain retains, memory; and how information is arranged while speaking, expression. Many frightening claims, when unpacked, land only on the lightest, third layer.

Behind these questions is linguist Dan Slobin's precise phrase: \textbf{thinking for speaking}. Language trains not your worldview itself, but habits of organizing attention at the moment of speech. Turkish speakers maintain professional-level sensitivity to ``Who said this? How was it learned?'' because every utterance requires taking stock. But that sensitivity follows the language switch; it is not engraved on the retina.

A familiar example secures the distinction. Years on city streets give taxi drivers astonishing mental maps; name an intersection, and they can report its direction. You would not say ``the taxi industry determines the only way they can see the world.'' You would say their occupation trained attention. \textbf{Training attention is not imprisoning thought.} This distinction is the chapter's backbone: language does much to guide attention and memory; the strong version claims it locks thought's possibilities. The gap between those is what science must measure.

Install one final alarm, because a mistake is especially tempting: observing that a grammar does not obligatorily mark the future, then inferring ``these people do not plan ahead.'' This mistakes ``grammar does not tax it'' for ``the mind has no room for it.'' By that logic, English does not obligatorily mark whether events happen before or after noon; can English speakers not distinguish morning from afternoon? A grammatical checklist is a tariff schedule, not the mind's complete furniture. Untaxed goods still fill the shelves.

\section{Why Is the Sea ``Wine-Colored''?}
Now read a little primary material from nearly three thousand years ago.

The ancient Greek Homeric epics repeatedly give the sea a fixed description. The epics include the Iliad and the Odyssey; the latter tells of Odysseus's ten-year sea journey home. Rendered literally, the sea is called again and again:

\begin{quote}
``The wine-colored sea.'' (From Homer's Odyssey; the same formula also recurs in the Iliad.)
\end{quote}
Pause. What color is wine? Dark red, deep purple, yellow-brown, perhaps, but not the ``azure sea'' that comes readily to our lips. Nor is this an isolated oddity: it is an epic formula, often recurring whenever the sea is sung.

The first prominent person to take serious issue with this was William Gladstone, later four times Britain's prime minister. Outside politics, he was a passionate Homer enthusiast. In 1858 he published the substantial Studies on Homer and the Homeric Age, doing laborious work: counting all the epics' color terms. The result startled him. Homer's palette was strange: wine-colored sea, iron's color applied to sky and water, almost no dedicated name for blue. Gladstone did not dodge the anomaly. He boldly guessed that ancient Greek eyes might differ from ours, their color discrimination not yet fully developed.

Today that guess is wrong; we will see evidence shortly. But the fact he counted is genuine and checkable: Homeric color vocabulary does not match our familiar palette.

Before pitying the ancient Greeks, look at Chinese. Qīng has a territory broad enough to overwhelm learners. ``Your Collar,'' in the Odes of Zheng section of the Book of Songs, says:

\begin{quote}
Deep qīng is your collar; long, long is my longing. (Book of Songs, Odes of Zheng, ``Your Collar'')
\end{quote}
Here qīng describes the collar's rich, dark color. The same character means blue in ``qīng sky,'' green in ``qīng mountains,'' black in ``qīng hair.'' One word spans three compartments of our palette.

Now ask about ourselves: have Chinese speakers across the centuries been unable to distinguish blue, green, and black? Of course not. Painters can, dyers can, and jade-buying aunties certainly can. Homer's sailors could distinguish wine from seawater too; failure would wreck ships. \textbf{Vocabulary's boundaries are not perception's boundaries.} A language's color grid records which distinctions it historically chose to tax, not which light its speakers' eyes lacked.

Why, then, could qīng's all-in-one pocket serve steadily for millennia? One plausible explanation is that its contents---vegetation, sky, deep water, hair---were grouped together in daily life as ``deep, cool colors.'' If life did not require separating them daily, language had no need for separate words. Ancient Greece was similar: blue as a tangible surface color or dye, something held or painted on a wall, was relatively uncommon in everyday objects. Sky and sea may be blue, but neither dyes cloth or fires pottery. Many languages name colors according to practical needs rather than spectral order. This belongs to ``how people then understood'': their use was practical. ``How we judge'' is separate. We need neither mock them for ``not seeing blue'' nor romanticize that ``their sea really was different.''

\section{One Child, Three Explanations}
Return to Queensland's red soil. We have an uncontested evidential fact: Guugu Yimithirr speakers, children included, have astonishing directional ability. Distinguish four things: what happens, their strong orientation; why, the arena of competing explanations; how participants understand it, seeing nothing remarkable because speech naturally includes directions; and how we evaluate it, perhaps calling it ``higher'' cognition. Set that value judgment aside. We compare why it happens.

\textbf{Explanation one: language trains it (language first).}

The reasoning: grammar supplies tens of thousands of compulsory daily exercises. From learning to speak, children mark directions in every sentence, effectively answering thousands of ``Which way is north now?'' questions. Ten years of practice would make poor orientation surprising. Late in the twentieth century, teams led by linguist Stephen Levinson ran comparative experiments in several communities using absolute directions. Asked to remember arrangements on a tabletop, these speakers remembered compass orientation; speakers of left/right languages remembered relative to their own bodies. Memory habits aligned closely with grammar. This account is neat and experimentally supported. Its limitation is that correlation does not establish causation: a third factor may have nurtured both language habits and spatial ability.

\textbf{Explanation two: environment and lifestyle train it (life first).}

The reasoning: gathering and hunting across open savanna makes orientation a survival skill. Finding water, navigating, and returning home depend on it. Language merely fixes this skill into grammar, as fishing communities develop many names for boat parts. Vocabulary does not create fishers; fishing creates vocabulary. Its strength is intuitive: skills grow around livelihoods. Its limitation is that environment explains the need for orientation, not why grammar specifically chooses north, south, east, and west as coordinates. Environment does not invent grammar; no jungle sign says ``Please use absolute directions.''

\textbf{Explanation three: language, environment, and culture shape one another (coevolution).}

The reasoning: perhaps the question itself is wrong. Language shapes attention, attention skills, skills lifestyles, and lifestyles language again. This loop has turned many times; nowhere on it yields a clear starting point. Savanna life needs orientation, so language encodes direction; grammar makes every child report directions, strengthening the skill and deepening the lifestyle. This account is the most comprehensive and closest to the real world's mess. Its limitation is just as clear: comprehensive explanations can be hard to test. Scientific progress needs predictions of ``If I am right, we should observe A rather than B.'' ``They influence one another'' fits almost any observation.

Each explanation holds part of the truth and falls short of the whole. This is not fence-sitting but a methodological reminder: \textbf{when an explanation explains everything, it has often discarded something in advance.} Remember this red soil when offered packaged answers that ``language determines everything'' or ``language determines nothing.''

\section{Further Challenge: Testing the Hypothesis}
Now for the chapter's weightiest stage: what has a century of science made of the Sapir--Whorf hypothesis?

First, the strong version's fate. Two solid nails were driven into its coffin.

The first was Whorf's own Hopi case. His claim that Hopi people lacked a concept of time astonished the world, yet for more than half a century few went to live in Hopi communities to check. Linguist Ekkehart Malotki eventually conducted years of fieldwork and published the substantial Hopi Time in 1983. It documents abundant temporal expressions: words for yesterday, tomorrow, periods of time, and grammatical means of ordering events. Hopi people naturally distinguish seasons and the schedules of weddings and funerals; a people unable to distinguish time could not grow corn. Whorf mistook grammatical packaging for a people's mental contents. This is our most important counterexample and easiest mistake: \textbf{inferring ``speakers lack a concept'' from ``their grammar lacks a category.''} That leap lands in a pit nine times out of ten. Guard against it, and you will stop using language structure to draw whole peoples' personality portraits.

The second nail was color. Gladstone guessed that ancient Greek eyes had not fully developed. Berlin and Kay, and much subsequent research, established that humanity shares basic color-discrimination abilities. Differences concern which distinctions languages choose to name. Genetically caused color blindness is another matter, belonging to eyes, not language. Homer's sailors looked at the same sea as us, the same light passing over their retinas.

The strong version fell. But inspect the weak version's experimental scores carefully: they are real, repeatable, and interesting.

\textbf{Color:} Russian has no general term matching English blue, instead treating siniy, ``dark blue,'' and goluboy, ``light blue,'' as separate basic colors, as naturally as we distinguish red and pink. A 2007 experiment in Proceedings of the National Academy of Sciences asked Russian and English speakers to distinguish blue patches. Russian speakers were faster when the patches crossed their dark/light-blue word boundary. Within the boundary, the advantage disappeared. Requiring them simultaneously to remember digits, occupying language resources, also weakened it. Notice the result: language participates in real-time perception, but changes \textbf{speed and habit}, not possibility. Given a little time, English speakers distinguish the blues too.

\textbf{Direction:} Levinson's experiments, as discussed, find that communities' memory for arrangements follows their grammars. This is evidence of linguistic habits entering memory.

\textbf{Time:} English speaks of time through front and back, looking forward toward the future. Chinese also uses up and down: ``upper week'' for last week, ``lower month'' for next month. Time can flow vertically in Chinese. More strikingly, research published in 2006 recorded temporal gestures among Aymara speakers in the South American Andes: the past is \textbf{in front}, because it has been seen and lies before the eyes; the future is \textbf{behind}, because it cannot be seen. For the same human species, time can flow in at least three directions.

\textbf{Number:} The Munduruku of the Brazilian Amazon have number words extending approximately to five. In 2004, a French research team reported in Science that exact calculation above five was difficult for them, while estimating ``roughly which pile is larger'' remained intact. This suggests that approximate quantity sense is factory hardware, while exact counting's ``machine'' really needs linguistic and symbolic scaffolding. A related point for you: Chinese eleven and twelve are transparently ``ten-one'' and ``ten-two,'' exposing place value. English eleven and twelve do not reveal that structure. Some research suggests this may be one reason East Asian children acquire tens-and-ones concepts somewhat earlier. Language's tariff schedule reaches even elementary mathematics.

Put the four domains together, and the verdict reads: \textbf{the strong version is largely falsified; the weak version holds, but within careful limits}. Language influences attention, memory, and speed at the level of habits, not thought's possibilities. Your native language is a gym, not a prison.

What about people with two languages? Bilinguals are a natural laboratory. Research and innumerable personal reports suggest the same picture: they are not two brains taking turns. Describe the same film in different languages, and the same people's emphases genuinely shift. Many say ``I love you'' feels weightier, insults more painful, in their native language. Aneta Pavlenko's years studying multilingual emotional language record many such reports: later-learned languages often put emotion behind a pane of glass. Other research finds that the same people give more ``cool-headed'' answers to moral dilemmas in a foreign language. Linguistic distance itself can be a resource.

Keep the right mental model: \textbf{bilinguals do not possess two worldviews, but a larger toolbox}. Switching language changes lenses, not eyes. This is among the deepest benefits of learning another language: you first discover that ``natural'' native-language distinctions---such as ``have eaten/am eating'' marking actions in time---are choices, not necessities. A foreign language mirrors your native one, revealing your own previously invisible contours.

Finally, draw the red line boldly. None of these experiments supports ``speaking language X means having personality Y.'' They concern what a grammar habitually highlights, not what kind of people a nation contains. Any inference sliding from ``their language lacks a word'' to ``they are shortsighted, cold, or careless'' turns statistical habits into individual destinies. Historically, that slide often wears prejudice's clothing and serves contempt for others. You have seen its mechanism; keep immunity for life.

\section{Today: New Words, Translation, and Two Blues}
This old question plays out in your phone every day.

Consider new words. Nèijuǎn, ``involution''; tǎngpíng, ``lying flat''; pòfáng, ``having one's defenses breached'' emotionally; emo. Chinese sprouts a new crop regularly. New words act like flashlights: once one appears, a vague, unsayable feeling is illuminated, named, and discussable. Before involution became popular, ``everyone keeps raising the stakes but nobody fares better'' certainly existed. Without a name it was harder to discuss publicly. This demonstrates the weak version: words issue thought a visa, not a birth certificate. Yet the same flashlight can restrict vision to its beam. Call every conflict ``PUA,'' manipulative psychological control, and real manipulation blurs with ordinary friction. Words mark running lanes for thought; what lies outside can become a blind spot.

Now translation. Chinese has five familiar distinctions: bóbo, father's elder brother; shūshu, father's younger brother; jiùjiu, mother's brother; gūfu, father's sister's husband; yífu, mother's sister's husband. English calls them all uncle. English people do not lack relatives; their language simply does not tax this distinction. Conversely, English privacy long lacked a convenient Chinese equivalent; yǐnsī gained currency later. Many concepts now seeming obvious entered Chinese's tariff schedule only recently. Every language is an attention tariff; learning one means reading another civilization's tax history.

People keep modifying language too. Over a century ago, written Chinese's third-person pronoun did not distinguish men and women: both used the same tā character. During the New Culture Movement, Liu Bannong and others promoted a different character, also pronounced tā, specifically for women. It existed in old books but was extremely rare; this campaign revived it. In 1920 Liu wrote the famous poem ``How Can I Help Thinking of Her?'' A century later, young internet users invented uppercase TA to bypass the gender binary. Some changes remain; others disappear. The question was never whether humans can remodel language; they always have. It is whether thought and behavior follow linguistic changes. You will witness several rounds of testing during your lifetime.

Then there is you. If you have learned some English and inhabit a Chinese-English online mix, you already live a miniature bilingual life. Next time you naturally say a Chinese sentence containing deadline, pause half a second and see what you did. Your Chinese was not ``polluted.'' You selected the handiest tool from a larger box.

\section{Over to You: Requiring an Evidence Marker}
End with a thought experiment.

Imagine a language reform: from today, Chinese adds a rule. Every factual statement requires a small verb marker declaring its source: one for ``I saw it,'' one for ``I heard it,'' another for ``I infer it.'' Turkish and Amazonian Tuyuca genuinely have such machinery; speakers cannot simply omit it.

Install it in today's world: class-chat rumors, viral trends, conversation screenshots, ``a friend told me.'' Before each utterance, grammar makes the speaker inventory their evidence.

\textbf{Would this society be more honest?}

Give your answer and reasons. Before concluding, count your resources again.

For reform: as tense makes us count time, this would make everyone count evidence. ``I heard'' would become a default setting rather than a courageous admission. Rumor-making would grow costlier because every retelling must leave fingerprints.

Against reform: it cannot cure intentional lying. Source marking and honesty differ; deceivers can learn to mark ``I saw it'' without blinking. It may also give eyewitness claims undeserved authority, though eyes are deceivable. More deeply, the whole chapter reminds us that language can reshape habits of attention, but grammar cannot single-handedly reshape the motives for attending.

Go deeper still: if compulsory sourcing really produced a society-wide habit of checking evidence, did grammar transform thought, or was thought waiting for an excuse? In other words, \textbf{when a language changes, does it push society forward, or has society already lifted its foot and merely borrowed language's hand to set it down?}

There is no standard answer. But from now on, whenever you say ``As everyone knows,'' you may hear this chapter quietly ask: who knows? How? Did you hear it yourself?
